\section{Step Responses}

The closed-loop step responses were generated at the nominal sustainer point:

\[
    X_{CG}=2.6012~\text{m}, \quad
    m=556.921~\text{kg}, \quad
    I_y=1611.097~\text{kg.m}^2.
\]

\subsection{Attitude Response}

The attitude loop tracks a unit attitude command with small overshoot and unit steady-state gain. The response is initially fast, and the remaining slow component is caused by the lower-frequency closed-loop pole.

\begin{figure}[H]
\centering
\includegraphics[width=0.88\textwidth]{figures/step_attitude.png}
\caption{Closed-loop attitude response to a unit step.}
\end{figure}

\subsection{Acceleration Response}

The acceleration loop was tested both with a unit command and with a \(5g\) command. The response has noticeable overshoot, but it remains stable and converges to the commanded acceleration.

\begin{figure}[H]
\centering
\includegraphics[width=0.88\textwidth]{figures/step_acceleration_unit.png}
\caption{Closed-loop acceleration response to a unit step.}
\end{figure}

\begin{figure}[H]
\centering
\includegraphics[width=0.88\textwidth]{figures/step_acceleration_5g.png}
\caption{Closed-loop acceleration response to a \(5g\) command.}
\end{figure}

For the nominal point, the MATLAB step metrics were:

\begin{table}[H]
\centering
\caption{Nominal closed-loop step metrics.}
\begin{tabular}{lrrrr}
\toprule
Autopilot & Rise time (s) & Settling time (s) & Overshoot (\%) & DC gain \\
\midrule
Attitude & 0.228 & 2.562 & 0.0 & 1.000 \\
Acceleration & 0.564 & 6.218 & 43.7 & 1.000 \\
\bottomrule
\end{tabular}
\end{table}
